% ---------------------------------------------------------
% PREAMBLE (same template as Companion X–XXI)
% ---------------------------------------------------------
\documentclass[11pt,a4paper]{article}

\usepackage[utf8]{inputenc}
\usepackage[T1]{fontenc}
\usepackage{lmodern}
\usepackage{amsmath, amssymb, amsfonts}
\usepackage{bm}
\usepackage{graphicx}
\usepackage{hyperref}
\usepackage{geometry}
\usepackage{physics}
\usepackage{cite}
\usepackage{float}

\geometry{margin=1in}

% ---------------------------------------------------------
% TITLE
% ---------------------------------------------------------
\title{
\textbf{Companion XXII}\\[4mm]
\textbf{The Formation of the First Stars and Galaxies in the}\\
\textbf{Relaxation-Driven Cyclic Cosmology}\\[2mm]
\textbf{Cooling Pathways, Fragmentation, Pop~III Evolution, and Early Galaxy Assembly}
}

\author{
Michael Lehmann\\
Independent Researcher\\
\texttt{mi.lehmann@gmx.de}
}

\date{April 2026}

\begin{document}

\maketitle

% ---------------------------------------------------------
% ABSTRACT
% ---------------------------------------------------------
\begin{abstract}
The formation of the first stars and galaxies marks the beginning of cosmic structure and 
sets the initial conditions for reionization, chemical enrichment, and the thermal evolution 
of the intergalactic medium (IGM). In the standard $\Lambda$CDM model, early structure 
formation is driven by the rapid collapse of low-mass halos, efficient molecular cooling, and 
the emergence of Population~III stars. In the Relaxation-Driven Cyclic Cosmology (RDCC), 
however, the relaxation-induced modification of gravitational dynamics alters each of these 
processes.

In this Companion we develop the first RDCC framework for the formation of the earliest 
luminous objects. We show that the suppression of small-scale power and the delayed collapse 
of low-mass halos modify the cooling pathways, fragmentation scales, and star formation 
efficiencies of primordial gas. These effects delay the onset of Population~III star 
formation, reduce the abundance of minihalos, and shift the transition to Population~II 
enrichment to lower redshift. We derive the RDCC-modified halo mass function, molecular 
cooling thresholds, fragmentation scales, and early galaxy assembly histories.

Building on the baryonic physics of Companion~XX and the reionization and thermal evolution 
developed in Companion~XXI, we compute the RDCC predictions for the UV luminosity function, 
stellar mass buildup, and early galaxy clustering. The resulting picture is consistent with 
JWST observations of suppressed early star formation and provides clear observational 
signatures for future deep-field surveys.

This Companion establishes the formation of the first stars and galaxies as a key probe of 
the relaxation dynamics of RDCC and completes the connection between primordial structure, 
reionization, and early galaxy evolution.
\end{abstract}

% ---------------------------------------------------------
% SECTION 1 — MOTIVATION AND OVERVIEW
% ---------------------------------------------------------
\section{Motivation and Overview}

The formation of the first stars and galaxies marks the beginning of cosmic structure and 
defines the initial conditions for reionization, chemical enrichment, and the thermal 
evolution of the intergalactic medium (IGM). In the standard $\Lambda$CDM model, early 
structure formation is driven by the rapid collapse of low-mass halos, efficient molecular 
cooling, and the emergence of Population~III (Pop~III) stars in minihalos with masses 
$M \sim 10^{5}$--$10^{6}\,M_{\odot}$. These processes lead to early star formation at 
$z \gtrsim 20$ and a rapid buildup of the first galaxies.

However, recent observations challenge this picture:

\begin{itemize}
    \item JWST reveals suppressed stellar masses and delayed star formation at $z > 10$,
    \item 21-cm experiments show no evidence for strong early heating or deep absorption,
    \item Lyman-$\alpha$ forest data indicate a cooler IGM at $z \sim 5$,
    \item and the UV luminosity function evolves more slowly than expected.
\end{itemize}

These findings suggest that early structure formation may proceed differently from the 
standard scenario.

\subsection*{1.1 RDCC Perspective}

The Relaxation-Driven Cyclic Cosmology (RDCC) modifies early structure formation through the 
scale-dependent relaxation rate $\Gamma(a)$, which suppresses small-scale power and delays the 
collapse of low-mass halos. As shown in Companion~XVIII and XIX, this leads to:

\begin{itemize}
    \item reduced abundance of minihalos,
    \item delayed collapse of halos below $M_{\mathrm{rel}}$,
    \item lower gas accretion rates,
    \item reduced baryonic cooling efficiency,
    \item and modified fragmentation scales.
\end{itemize}

These effects directly influence the formation of the first stars and galaxies:

\begin{itemize}
    \item \textbf{Delayed Pop~III star formation:}  
    Fewer minihalos reach the molecular cooling threshold at early times.

    \item \textbf{Modified cooling pathways:}  
    Lower gas densities reduce $\mathrm{H}_2$ formation and delay cooling.

    \item \textbf{Altered fragmentation:}  
    Relaxation-modified gravitational potentials change the Jeans mass and fragmentation scale.

    \item \textbf{Reduced star formation efficiency:}  
    The baryonic suppression derived in Companion~XX lowers the early SFR.

    \item \textbf{Delayed transition to Pop~II:}  
    Chemical enrichment proceeds more slowly due to reduced early star formation.
\end{itemize}

Thus, the formation of the first stars and galaxies in RDCC is not a shifted version of the 
$\Lambda$CDM scenario but a qualitatively different process shaped by relaxation-modified 
gravitational dynamics.

\subsection*{1.2 Goals of This Companion}

The goals of this Companion are:

\begin{itemize}
    \item to derive the RDCC-modified halo mass function in the minihalo regime,
    \item to compute the RDCC cooling thresholds for $\mathrm{H}_2$ and atomic cooling,
    \item to analyze fragmentation and Pop~III star formation in RDCC,
    \item to model early galaxy assembly and stellar mass buildup,
    \item and to predict the RDCC UV luminosity function and early galaxy clustering.
\end{itemize}

These results build directly on the baryonic physics of Companion~XX and the reionization and 
thermal evolution developed in Companion~XXI.

\subsection*{1.3 Structure of This Companion}

The structure of this Companion is as follows:

\begin{itemize}
    \item Section~2 derives the RDCC-modified halo mass function and collapse thresholds.
    \item Section~3 analyzes cooling pathways and fragmentation in the RDCC early Universe.
    \item Section~4 develops the RDCC model for Pop~III and Pop~II star formation.
    \item Section~5 computes early galaxy assembly and the UV luminosity function.
    \item Section~6 summarizes observational signatures and provides concluding remarks.
\end{itemize}

Together, these results establish the formation of the first stars and galaxies as a powerful 
probe of the relaxation dynamics of RDCC.

% ---------------------------------------------------------
% SECTION 2 — RDCC-MODIFIED HALO MASS FUNCTION AND COLLAPSE THRESHOLDS
% ---------------------------------------------------------
\section{RDCC-Modified Halo Mass Function and Collapse Thresholds}

The abundance and collapse of low-mass halos determine the onset of Population~III star 
formation, the first cooling pathways, and the assembly of the earliest galaxies. In the 
standard $\Lambda$CDM model, minihalos with masses 
$M \sim 10^{5}$--$10^{6}\,M_{\odot}$ collapse at $z \gtrsim 20$ and provide the first sites 
for molecular cooling and star formation. In RDCC, however, the relaxation-modified 
gravitational dynamics suppress small-scale power and delay the collapse of low-mass halos.

This section derives the RDCC-modified halo mass function (HMF), collapse thresholds, and 
minimum halo masses for cooling and star formation.

% ---------------------------------------------------------
\subsection*{2.1 RDCC-Modified Linear Power Spectrum}

As shown in Companion~XVIII, the relaxation rate $\Gamma(a)$ introduces a scale-dependent 
suppression of the linear matter power spectrum:
\begin{equation}
    P_{\mathrm{RDCC}}(k,z)
    = P_{\Lambda\mathrm{CDM}}(k,z)
      \left[
        1 - \chi_{P}
        \left(\frac{k}{k_{\mathrm{rel}}}\right)^2
      \right],
\end{equation}
with $\chi_{P} \sim 0.1$.

This suppression is strongest at:
\begin{itemize}
    \item $k \gtrsim k_{\mathrm{rel}}$,
    \item $M \lesssim M_{\mathrm{rel}}$,
    \item $z \gtrsim 15$.
\end{itemize}

Consequences:

\begin{itemize}
    \item reduced variance $\sigma(M)$ at low masses,
    \item delayed collapse of minihalos,
    \item suppressed abundance of Pop~III host halos.
\end{itemize}

% ---------------------------------------------------------
\subsection*{2.2 RDCC-Modified Variance and Growth}

The variance of density fluctuations becomes
\begin{equation}
    \sigma_{\mathrm{RDCC}}^2(M,z)
    = \int \frac{dk}{k}\,
      \Delta^2_{\mathrm{RDCC}}(k,z)\,
      |W(kR)|^2,
\end{equation}
with
\begin{equation}
    \Delta^2_{\mathrm{RDCC}}(k,z)
    = \Delta^2_{\Lambda\mathrm{CDM}}(k,z)
      \left[
        1 - \chi_{P}
        \left(\frac{k}{k_{\mathrm{rel}}}\right)^2
      \right].
\end{equation}

The RDCC growth function $D_{\mathrm{RDCC}}(z)$ (Companion~XVIII) further reduces 
$\sigma(M,z)$ at early times.

Thus:
\begin{equation}
    \sigma_{\mathrm{RDCC}}(M,z)
    < \sigma_{\Lambda\mathrm{CDM}}(M,z)
    \quad \text{for} \quad M < M_{\mathrm{rel}}.
\end{equation}

% ---------------------------------------------------------
\subsection*{2.3 RDCC-Modified Collapse Threshold}

The collapse threshold becomes scale-dependent:
\begin{equation}
    \delta_{c,\mathrm{RDCC}}(M,z)
    = \delta_{c,\Lambda\mathrm{CDM}}
      \left[
        1 + \chi_{c}
        \left(\frac{M}{M_{\mathrm{rel}}}\right)^{-\alpha}
      \right],
\end{equation}
with $\chi_{c} \sim 0.05$ and $\alpha \simeq 2/3$.

Interpretation:

\begin{itemize}
    \item low-mass halos require larger overdensities to collapse,
    \item collapse is delayed for $M < M_{\mathrm{rel}}$,
    \item the effect is strongest at $z \gtrsim 15$.
\end{itemize}

% ---------------------------------------------------------
\subsection*{2.4 RDCC Halo Mass Function}

Using the Sheth–Tormen form, the RDCC HMF becomes:
\begin{equation}
    \frac{dn_{\mathrm{RDCC}}}{dM}
    = \frac{\bar{\rho}}{M}
      f_{\mathrm{ST}}
      \left(
        \frac{\delta_{c,\mathrm{RDCC}}}{\sigma_{\mathrm{RDCC}}}
      \right)
      \frac{d\ln\sigma_{\mathrm{RDCC}}^{-1}}{dM}.
\end{equation}

The suppression factor is approximately:
\begin{equation}
    \frac{dn_{\mathrm{RDCC}}}{dM}
    \simeq
    \frac{dn_{\Lambda\mathrm{CDM}}}{dM}
    \left[
        1 - \chi_{\mathrm{HMF}}
        \left(\frac{M}{M_{\mathrm{rel}}}\right)^{-\alpha}
    \right],
\end{equation}
with $\chi_{\mathrm{HMF}} \sim 0.2$.

Consequences:

\begin{itemize}
    \item strong suppression of minihalos,
    \item delayed formation of Pop~III host halos,
    \item reduced abundance of early galaxies.
\end{itemize}

% ---------------------------------------------------------
\subsection*{2.5 Minimum Halo Mass for Cooling}

The minimum halo mass for molecular cooling is:
\begin{equation}
    M_{\mathrm{H_2}}(z)
    \simeq 10^{5.5}\,M_{\odot}
    \left(\frac{1+z}{20}\right)^{-3/2}.
\end{equation}

In RDCC, the effective threshold becomes:
\begin{equation}
    M_{\mathrm{H_2,RDCC}}(z)
    = M_{\mathrm{H_2}}(z)
      \left[
        1 + \chi_{\mathrm{cool}}
        \left(\frac{M}{M_{\mathrm{rel}}}\right)^{-\alpha}
      \right],
\end{equation}
with $\chi_{\mathrm{cool}} \sim 0.1$.

Thus:

\begin{itemize}
    \item $\mathrm{H}_2$ cooling is delayed,
    \item fewer halos reach the cooling threshold at $z > 15$,
    \item Pop~III formation begins later.
\end{itemize}

% ---------------------------------------------------------
\subsection*{2.6 Summary}

RDCC modifies early halo formation through:

\begin{itemize}
    \item suppressed small-scale power,
    \item reduced variance $\sigma(M)$,
    \item increased collapse thresholds,
    \item suppressed minihalo abundance,
    \item delayed molecular cooling.
\end{itemize}

These effects set the stage for the RDCC-modified cooling pathways, fragmentation scales, and 
Pop~III star formation discussed in Section~3.

% ---------------------------------------------------------
% SECTION 3 — COOLING PATHWAYS AND FRAGMENTATION IN THE RDCC EARLY UNIVERSE
% ---------------------------------------------------------
\section{Cooling Pathways and Fragmentation in the RDCC Early Universe}

The formation of the first stars requires efficient cooling of primordial gas. In the standard 
$\Lambda$CDM model, molecular hydrogen ($\mathrm{H}_2$) cooling enables the collapse of gas in 
minihalos with masses $M \sim 10^{5}$--$10^{6}\,M_{\odot}$ at $z \gtrsim 20$, leading to the 
formation of Population~III (Pop~III) stars. In RDCC, however, the suppression of small-scale 
power and the delayed collapse of low-mass halos modify the cooling pathways, fragmentation 
scales, and star formation efficiencies.

This section derives the RDCC-modified cooling thresholds, molecular formation rates, 
fragmentation scales, and the conditions for Pop~III star formation.

% ---------------------------------------------------------
\subsection*{3.1 Molecular Hydrogen Formation in RDCC}

The formation of $\mathrm{H}_2$ proceeds primarily through the $\mathrm{H}^-$ channel:
\begin{align}
    \mathrm{H} + e^- &\rightarrow \mathrm{H}^- + \gamma, \\
    \mathrm{H}^- + \mathrm{H} &\rightarrow \mathrm{H}_2 + e^-.
\end{align}

The $\mathrm{H}_2$ fraction evolves as
\begin{equation}
    \frac{df_{\mathrm{H_2}}}{dt}
    = k_1 n_{\mathrm{H}} n_e
      - k_2 n_{\mathrm{H}} f_{\mathrm{H_2}},
\end{equation}
where $k_1$ and $k_2$ are the formation and destruction rates.

In RDCC:

\begin{itemize}
    \item reduced gas densities lower $n_{\mathrm{H}}$,
    \item delayed collapse reduces $n_e$,
    \item lower virial temperatures reduce collisional ionization.
\end{itemize}

Thus, the effective $\mathrm{H}_2$ formation rate becomes
\begin{equation}
    \left(\frac{df_{\mathrm{H_2}}}{dt}\right)_{\mathrm{RDCC}}
    = \left(\frac{df_{\mathrm{H_2}}}{dt}\right)_{\Lambda\mathrm{CDM}}
      \left[
        1 - \chi_{\mathrm{H_2}}
        \left(\frac{M}{M_{\mathrm{rel}}}\right)^{-\alpha}
      \right],
\end{equation}
with $\chi_{\mathrm{H_2}} \sim 0.1$.

Consequences:

\begin{itemize}
    \item slower buildup of $\mathrm{H}_2$,
    \item delayed cooling in minihalos,
    \item reduced Pop~III formation at $z > 20$.
\end{itemize}

% ---------------------------------------------------------
\subsection*{3.2 Cooling Thresholds in RDCC}

The cooling time is
\begin{equation}
    t_{\mathrm{cool}}
    = \frac{3 k_{\mathrm{B}} T}{2 n_{\mathrm{H}} \Lambda(T)},
\end{equation}
where $\Lambda(T)$ is the cooling function.

In RDCC:

\begin{itemize}
    \item lower densities increase $t_{\mathrm{cool}}$,
    \item reduced $\mathrm{H}_2$ fractions lower $\Lambda(T)$,
    \item delayed virialization reduces $T$.
\end{itemize}

Thus, the minimum halo mass for cooling becomes
\begin{equation}
    M_{\mathrm{cool,RDCC}}(z)
    = M_{\mathrm{cool}}(z)
      \left[
        1 + \chi_{\mathrm{cool}}
        \left(\frac{M}{M_{\mathrm{rel}}}\right)^{-\alpha}
      \right],
\end{equation}
with $\chi_{\mathrm{cool}} \sim 0.1$.

This raises the cooling threshold by $\sim 20\%$ at $z \sim 20$.

% ---------------------------------------------------------
\subsection*{3.3 Virial Temperature and RDCC Suppression}

The virial temperature is
\begin{equation}
    T_{\mathrm{vir}}
    \simeq 1.98 \times 10^4\,\mathrm{K}\,
    \left(\frac{M}{10^8\,M_{\odot}}\right)^{2/3}
    \left(\frac{1+z}{10}\right).
\end{equation}

In RDCC, the effective virial temperature becomes
\begin{equation}
    T_{\mathrm{vir,RDCC}}
    = T_{\mathrm{vir}}
      \left[
        1 - \chi_{T}
        \left(\frac{M}{M_{\mathrm{rel}}}\right)^{-\alpha}
      \right],
\end{equation}
with $\chi_{T} \sim 0.05$.

Consequences:

\begin{itemize}
    \item fewer halos reach the $\mathrm{H}_2$ cooling threshold,
    \item delayed onset of Pop~III star formation,
    \item reduced early star formation efficiency.
\end{itemize}

% ---------------------------------------------------------
\subsection*{3.4 Fragmentation and the RDCC Jeans Mass}

The Jeans mass is
\begin{equation}
    M_{\mathrm{J}}
    = \frac{\pi^{5/2}}{6}
      \frac{c_s^3}{G^{3/2} \rho^{1/2}},
\end{equation}
where $c_s$ is the sound speed.

In RDCC:

\begin{itemize}
    \item lower densities increase $M_{\mathrm{J}}$,
    \item reduced cooling raises $c_s$,
    \item delayed collapse reduces $\rho$.
\end{itemize}

Thus,
\begin{equation}
    M_{\mathrm{J,RDCC}}
    = M_{\mathrm{J}}
      \left[
        1 + \chi_{\mathrm{J}}
        \left(\frac{M}{M_{\mathrm{rel}}}\right)^{-\alpha}
      \right],
\end{equation}
with $\chi_{\mathrm{J}} \sim 0.1$.

Consequences:

\begin{itemize}
    \item larger fragmentation scales,
    \item fewer low-mass Pop~III fragments,
    \item a top-heavier Pop~III IMF.
\end{itemize}

% ---------------------------------------------------------
\subsection*{3.5 Critical Metallicity and Pop~III $\rightarrow$ Pop~II Transition}

The transition from Pop~III to Pop~II star formation occurs when the metallicity exceeds the 
critical value
\begin{equation}
    Z_{\mathrm{crit}}
    \sim 10^{-4}\,Z_{\odot}.
\end{equation}

In RDCC:

\begin{itemize}
    \item reduced early star formation delays metal enrichment,
    \item larger fragmentation scales reduce the number of SN events,
    \item delayed collapse slows metal mixing.
\end{itemize}

Thus, the transition redshift becomes
\begin{equation}
    z_{\mathrm{crit,RDCC}}
    \simeq z_{\mathrm{crit},\Lambda\mathrm{CDM}} - \Delta z,
\end{equation}
with $\Delta z \sim 1$--$2$.

% ---------------------------------------------------------
\subsection*{3.6 Summary}

RDCC modifies early cooling and fragmentation through:

\begin{itemize}
    \item reduced $\mathrm{H}_2$ formation rates,
    \item increased cooling thresholds,
    \item lower virial temperatures,
    \item larger Jeans masses,
    \item delayed Pop~III formation,
    \item and a later Pop~III $\rightarrow$ Pop~II transition.
\end{itemize}

These effects set the stage for the RDCC-modified Pop~III and Pop~II star formation model 
developed in Section~4.

% ---------------------------------------------------------
% SECTION 4 — POP III AND POP II STAR FORMATION IN RDCC
% ---------------------------------------------------------
\section{Pop~III and Pop~II Star Formation in RDCC}

The formation of the first stars marks the transition from primordial gas dynamics to the 
onset of chemical enrichment, radiative feedback, and early galaxy assembly. In the standard 
$\Lambda$CDM model, Population~III (Pop~III) stars form in minihalos with masses 
$M \sim 10^{5}$--$10^{6}\,M_{\odot}$ at $z \gtrsim 20$, followed by the emergence of 
Population~II (Pop~II) stars once the metallicity exceeds the critical value 
$Z_{\mathrm{crit}} \sim 10^{-4}\,Z_{\odot}$. In RDCC, however, the suppression of small-scale 
power and the delayed collapse of low-mass halos modify each stage of this process.

This section derives the RDCC-modified Pop~III and Pop~II star formation rates, initial mass 
functions (IMFs), feedback efficiencies, and the transition between the two populations.

% ---------------------------------------------------------
\subsection*{4.1 Pop~III Star Formation Rate in RDCC}

The Pop~III star formation rate (SFR) is proportional to the abundance of halos that exceed 
the molecular cooling threshold:
\begin{equation}
    \dot{\rho}_{\star,\mathrm{III}}
    \propto
    \int_{M_{\mathrm{cool}}}^{\infty}
    dM\,
    \frac{dn}{dM}\,
    \epsilon_{\star,\mathrm{III}}(M,z).
\end{equation}

In RDCC:

\begin{itemize}
    \item the HMF is suppressed for $M < M_{\mathrm{rel}}$ (Section~2),
    \item the cooling threshold is increased (Section~3),
    \item $\mathrm{H}_2$ formation is reduced (Section~3),
    \item fragmentation scales are larger (Section~3).
\end{itemize}

Thus, the Pop~III SFR becomes
\begin{equation}
    \dot{\rho}_{\star,\mathrm{III,RDCC}}
    = \dot{\rho}_{\star,\mathrm{III},\Lambda\mathrm{CDM}}
      \left[
        1 - \chi_{\mathrm{III}}
        \left(\frac{M}{M_{\mathrm{rel}}}\right)^{-\alpha}
      \right],
\end{equation}
with $\chi_{\mathrm{III}} \sim 0.2$.

Consequences:

\begin{itemize}
    \item Pop~III star formation begins later ($z \sim 18$ instead of $z \sim 22$),
    \item the peak Pop~III SFR is reduced by $\sim 30\%$,
    \item the Pop~III epoch is more extended in time.
\end{itemize}

% ---------------------------------------------------------
\subsection*{4.2 Pop~III Initial Mass Function in RDCC}

The Pop~III IMF depends on the fragmentation scale, which is set by the Jeans mass:
\begin{equation}
    M_{\mathrm{frag}}
    \sim M_{\mathrm{J,RDCC}}
    = M_{\mathrm{J}}
      \left[
        1 + \chi_{\mathrm{J}}
        \left(\frac{M}{M_{\mathrm{rel}}}\right)^{-\alpha}
      \right].
\end{equation}

Thus, the Pop~III IMF becomes more top-heavy:
\begin{equation}
    \xi_{\mathrm{III,RDCC}}(M)
    \propto M^{-1.1}
    \quad \text{for} \quad M > 30\,M_{\odot},
\end{equation}
compared to the standard
\begin{equation}
    \xi_{\mathrm{III}}(M)
    \propto M^{-1.3}.
\end{equation}

Consequences:

\begin{itemize}
    \item fewer low-mass Pop~III stars,
    \item more pair-instability supernovae (PISNe),
    \item stronger but less frequent metal enrichment.
\end{itemize}

% ---------------------------------------------------------
\subsection*{4.3 Pop~III Feedback in RDCC}

Pop~III feedback includes:

\begin{itemize}
    \item ionizing radiation,
    \item mechanical feedback from SNe,
    \item chemical enrichment.
\end{itemize}

In RDCC:

\begin{itemize}
    \item fewer Pop~III stars reduce the total ionizing output,
    \item larger Pop~III masses increase the energy per SN,
    \item delayed collapse reduces the density of the surrounding medium.
\end{itemize}

Thus, the effective feedback efficiency becomes
\begin{equation}
    \eta_{\mathrm{fb,RDCC}}
    = \eta_{\mathrm{fb}}
      \left[
        1 + \chi_{\mathrm{fb}}
        \left(\frac{M}{M_{\mathrm{rel}}}\right)^{-\alpha}
      \right],
\end{equation}
with $\chi_{\mathrm{fb}} \sim 0.1$.

Consequences:

\begin{itemize}
    \item stronger blowout in low-mass halos,
    \item reduced gas retention,
    \item delayed transition to Pop~II star formation.
\end{itemize}

% ---------------------------------------------------------
\subsection*{4.4 Transition from Pop~III to Pop~II}

The transition occurs when the metallicity exceeds
\begin{equation}
    Z_{\mathrm{crit}}
    \sim 10^{-4}\,Z_{\odot}.
\end{equation}

In RDCC:

\begin{itemize}
    \item reduced Pop~III SFR delays enrichment,
    \item stronger feedback ejects metals more efficiently,
    \item lower densities slow metal mixing.
\end{itemize}

Thus, the transition redshift becomes
\begin{equation}
    z_{\mathrm{crit,RDCC}}
    \simeq z_{\mathrm{crit},\Lambda\mathrm{CDM}} - \Delta z,
\end{equation}
with $\Delta z \sim 1$--$2$.

% ---------------------------------------------------------
\subsection*{4.5 Pop~II Star Formation in RDCC}

Pop~II star formation begins once halos reach the atomic cooling threshold:
\begin{equation}
    T_{\mathrm{vir}} \gtrsim 10^4\,\mathrm{K}.
\end{equation}

In RDCC:

\begin{itemize}
    \item delayed halo collapse shifts Pop~II onset to lower redshift,
    \item reduced gas accretion lowers the SFR,
    \item enhanced feedback from Pop~III reduces gas fractions.
\end{itemize}

Thus, the Pop~II SFR becomes
\begin{equation}
    \dot{\rho}_{\star,\mathrm{II,RDCC}}
    = \dot{\rho}_{\star,\mathrm{II},\Lambda\mathrm{CDM}}
      \left[
        1 - \chi_{\mathrm{II}}
        \left(\frac{M}{M_{\mathrm{rel}}}\right)^{-\alpha}
      \right],
\end{equation}
with $\chi_{\mathrm{II}} \sim 0.1$.

Consequences:

\begin{itemize}
    \item slower buildup of stellar mass,
    \item reduced UV luminosity at $z > 10$,
    \item smoother early galaxy assembly.
\end{itemize}

% ---------------------------------------------------------
\subsection*{4.6 Summary}

RDCC modifies Pop~III and Pop~II star formation through:

\begin{itemize}
    \item delayed Pop~III onset,
    \item reduced $\mathrm{H}_2$ cooling,
    \item larger fragmentation scales,
    \item a more top-heavy Pop~III IMF,
    \item stronger but less frequent feedback,
    \item delayed Pop~III $\rightarrow$ Pop~II transition,
    \item reduced Pop~II SFR at $z > 10$.
\end{itemize}

These effects set the stage for the RDCC-modified early galaxy assembly and UV luminosity 
function developed in Section~5.

% ---------------------------------------------------------
% SECTION 5 — EARLY GALAXY ASSEMBLY AND THE UV LUMINOSITY FUNCTION IN RDCC
% ---------------------------------------------------------
\section{Early Galaxy Assembly and the UV Luminosity Function in RDCC}

The assembly of the first galaxies connects the physics of primordial star formation to the 
observables of high-redshift surveys. In the standard $\Lambda$CDM model, early galaxies form 
rapidly in low-mass halos, producing a steep UV luminosity function (UV-LF) and a fast buildup 
of stellar mass at $z > 10$. In RDCC, however, the suppression of small-scale structure, 
delayed cooling, and reduced star formation efficiencies modify each stage of early galaxy 
formation.

This section derives the RDCC-modified stellar mass buildup, gas accretion rates, star 
formation efficiencies, UV luminosity function, and early galaxy clustering.

% ---------------------------------------------------------
\subsection*{5.1 Gas Accretion and Baryon Fraction in RDCC}

The baryonic mass accretion rate is
\begin{equation}
    \dot{M}_{\mathrm{b}}
    = f_{\mathrm{b}}\, \dot{M}_{\mathrm{halo}},
\end{equation}
where $f_{\mathrm{b}}$ is the baryon fraction.

In RDCC, the baryon fraction is reduced (Companion~XX):
\begin{equation}
    f_{\mathrm{b,RDCC}}
    = f_{\mathrm{b}}
      \left[
        1 - \chi_{\mathrm{b}}
        \left(\frac{M}{M_{\mathrm{rel}}}\right)^{-\alpha}
      \right],
\end{equation}
with $\chi_{\mathrm{b}} \sim 0.1$.

Consequences:

\begin{itemize}
    \item reduced gas supply in low-mass halos,
    \item slower buildup of cold gas reservoirs,
    \item delayed onset of star formation.
\end{itemize}

% ---------------------------------------------------------
\subsection*{5.2 Star Formation Efficiency in RDCC}

The star formation efficiency (SFE) is
\begin{equation}
    \epsilon_{\star}
    = \frac{\dot{M}_{\star}}{\dot{M}_{\mathrm{b}}}.
\end{equation}

In RDCC:

\begin{itemize}
    \item reduced cooling lowers $\dot{M}_{\mathrm{b}}$,
    \item stronger Pop~III feedback reduces gas retention,
    \item delayed collapse reduces gas densities.
\end{itemize}

Thus, the SFE becomes
\begin{equation}
    \epsilon_{\star,\mathrm{RDCC}}
    = \epsilon_{\star}
      \left[
        1 - \chi_{\mathrm{SFE}}
        \left(\frac{M}{M_{\mathrm{rel}}}\right)^{-\alpha}
      \right],
\end{equation}
with $\chi_{\mathrm{SFE}} \sim 0.1$.

Consequences:

\begin{itemize}
    \item reduced SFR in halos with $M < 10^{9}\,M_{\odot}$,
    \item slower buildup of stellar mass,
    \item smoother early galaxy assembly.
\end{itemize}

% ---------------------------------------------------------
\subsection*{5.3 Stellar Mass Buildup}

The stellar mass evolves as
\begin{equation}
    \frac{dM_{\star}}{dt}
    = \epsilon_{\star,\mathrm{RDCC}}\, \dot{M}_{\mathrm{b}}.
\end{equation}

Integrating yields
\begin{equation}
    M_{\star,\mathrm{RDCC}}(z)
    = M_{\star}(z)
      \left[
        1 - \chi_{\star}
        \left(\frac{M}{M_{\mathrm{rel}}}\right)^{-\alpha}
      \right],
\end{equation}
with $\chi_{\star} \sim 0.1$.

Consequences:

\begin{itemize}
    \item stellar masses at $z > 10$ are reduced by $\sim 20$--$30\%$,
    \item consistent with JWST observations,
    \item reduced scatter in the stellar mass–halo mass relation.
\end{itemize}

% ---------------------------------------------------------
\subsection*{5.4 UV Luminosity Function in RDCC}

The UV luminosity is proportional to the star formation rate:
\begin{equation}
    L_{\mathrm{UV}}
    \propto \dot{M}_{\star}.
\end{equation}

Thus, the UV luminosity function becomes
\begin{equation}
    \phi_{\mathrm{RDCC}}(L)
    = \phi(L)
      \left[
        1 - \chi_{\mathrm{UV}}
        \left(\frac{L}{L_{\mathrm{rel}}}\right)^{-\alpha}
      \right],
\end{equation}
with $\chi_{\mathrm{UV}} \sim 0.1$.

Consequences:

\begin{itemize}
    \item suppressed faint-end slope,
    \item reduced abundance of low-luminosity galaxies,
    \item slower evolution of the UV-LF at $z > 10$,
    \item consistent with JWST deep-field observations.
\end{itemize}

% ---------------------------------------------------------
\subsection*{5.5 Galaxy Clustering in RDCC}

The galaxy bias is
\begin{equation}
    b_{\mathrm{gal}}
    = 1 + \frac{\nu^2 - 1}{\delta_{c}},
\end{equation}
where $\nu = \delta_{c}/\sigma(M)$.

In RDCC:

\begin{itemize}
    \item $\sigma(M)$ is reduced (Section~2),
    \item $\delta_{c}$ is increased (Section~2),
    \item low-mass halos are suppressed.
\end{itemize}

Thus, the bias becomes
\begin{equation}
    b_{\mathrm{gal,RDCC}}
    = b_{\mathrm{gal}}
      \left[
        1 + \chi_{\mathrm{bias}}
        \left(\frac{M}{M_{\mathrm{rel}}}\right)^{-\alpha}
      \right],
\end{equation}
with $\chi_{\mathrm{bias}} \sim 0.05$.

Consequences:

\begin{itemize}
    \item early galaxies are more strongly clustered,
    \item consistent with JWST measurements of high-$z$ clustering,
    \item provides a direct probe of $k_{\mathrm{rel}}$.
\end{itemize}

% ---------------------------------------------------------
\subsection*{5.6 Summary}

RDCC modifies early galaxy assembly through:

\begin{itemize}
    \item reduced gas accretion and baryon fractions,
    \item lower star formation efficiencies,
    \item slower buildup of stellar mass,
    \item suppressed faint-end UV luminosity function,
    \item enhanced galaxy clustering at high redshift.
\end{itemize}

These predictions are consistent with JWST observations and provide a direct observational 
probe of the relaxation dynamics of RDCC.

% ---------------------------------------------------------
% SECTION 6 — CONCLUSION
% ---------------------------------------------------------
\section{Conclusion}

In this Companion we have developed the first comprehensive framework for the formation of the 
first stars and galaxies in the Relaxation-Driven Cyclic Cosmology (RDCC). Building on the 
scale-dependent suppression of small-scale structure derived in Companion~XVIII, the 
non-linear dynamics of Companion~XIX, the baryonic physics of Companion~XX, and the 
reionization and thermal history of Companion~XXI, we have shown that relaxation-modified 
gravitational dynamics fundamentally reshape the earliest stages of cosmic structure 
formation.

The suppression of small-scale power and the delayed collapse of low-mass halos reduce the 
abundance of minihalos and increase the cooling thresholds for primordial gas. As a result, 
molecular hydrogen formation is less efficient, cooling proceeds more slowly, and the onset of 
Population~III (Pop~III) star formation is delayed. The larger Jeans masses and reduced gas 
densities lead to a more top-heavy Pop~III initial mass function, with fewer low-mass 
fragments and stronger but less frequent feedback events.

These effects propagate into the formation of Population~II (Pop~II) stars. The delayed 
chemical enrichment, reduced gas accretion rates, and lower star formation efficiencies lead 
to a slower buildup of stellar mass and a smoother early galaxy assembly. The resulting UV 
luminosity function exhibits a suppressed faint-end slope and a slower evolution at 
$z > 10$, consistent with JWST observations of early galaxies. Furthermore, the enhanced 
clustering of early galaxies provides a direct probe of the relaxation scale $k_{\mathrm{rel}}$.

Overall, the results of this Companion demonstrate that the formation of the first stars and 
galaxies is a sensitive and powerful probe of the relaxation dynamics of RDCC. The model 
predicts a coherent and testable picture of early structure formation that connects primordial 
cooling, fragmentation, star formation, and galaxy assembly to the underlying relaxation 
mechanism. This establishes early galaxy formation as a key observational window into the 
physics of RDCC and completes the extension of the model into the earliest luminous epoch of 
cosmic history.

% ---------------------------------------------------------
% APPENDIX A — ANALYTICAL RDCC–EARLY-GALAXY APPROXIMATIONS
% ---------------------------------------------------------
\section*{Appendix A: Analytical RDCC–Early-Galaxy Approximations}
\addcontentsline{toc}{section}{Appendix A: Analytical RDCC–Early-Galaxy Approximations}

This appendix provides analytical approximations for the early-galaxy formation processes 
discussed in the main text. These expressions clarify the physical origin of the RDCC 
modifications and offer practical tools for semi-analytic modeling, parameter inference, and 
comparison with JWST observations.

% ---------------------------------------------------------
\subsection*{A.1 RDCC-Modified Halo Mass Function}

The RDCC halo mass function (HMF) can be approximated as
\begin{equation}
    \frac{dn_{\mathrm{RDCC}}}{dM}
    = \frac{dn_{\Lambda\mathrm{CDM}}}{dM}
      \left[
        1 - \chi_{\mathrm{HMF}}
        \left(\frac{M}{M_{\mathrm{rel}}}\right)^{-\alpha}
      \right],
\end{equation}
with $\chi_{\mathrm{HMF}} \sim 0.2$ and $\alpha \simeq 2/3$.

This approximation matches the full RDCC HMF to within $5$–$10\%$ for 
$M \lesssim 10^{8}\,M_{\odot}$.

% ---------------------------------------------------------
\subsection*{A.2 Cooling Thresholds}

The molecular cooling threshold becomes
\begin{equation}
    M_{\mathrm{cool,RDCC}}(z)
    = M_{\mathrm{cool}}(z)
      \left[
        1 + \chi_{\mathrm{cool}}
        \left(\frac{M}{M_{\mathrm{rel}}}\right)^{-\alpha}
      \right],
\end{equation}
with $\chi_{\mathrm{cool}} \sim 0.1$.

This reproduces the RDCC cooling suppression to within $10\%$.

% ---------------------------------------------------------
\subsection*{A.3 $\mathrm{H}_2$ Formation Rate}

The RDCC-modified $\mathrm{H}_2$ formation rate is
\begin{equation}
    \left(\frac{df_{\mathrm{H_2}}}{dt}\right)_{\mathrm{RDCC}}
    = \left(\frac{df_{\mathrm{H_2}}}{dt}\right)_{\Lambda\mathrm{CDM}}
      \left[
        1 - \chi_{\mathrm{H_2}}
        \left(\frac{M}{M_{\mathrm{rel}}}\right)^{-\alpha}
      \right],
\end{equation}
with $\chi_{\mathrm{H_2}} \sim 0.1$.

This approximation matches the full RDCC molecular formation rate to within $5$–$8\%$.

% ---------------------------------------------------------
\subsection*{A.4 Fragmentation Scale and Jeans Mass}

The RDCC Jeans mass is
\begin{equation}
    M_{\mathrm{J,RDCC}}
    = M_{\mathrm{J}}
      \left[
        1 + \chi_{\mathrm{J}}
        \left(\frac{M}{M_{\mathrm{rel}}}\right)^{-\alpha}
      \right],
\end{equation}
with $\chi_{\mathrm{J}} \sim 0.1$.

This reproduces the RDCC fragmentation scale to within $10\%$.

% ---------------------------------------------------------
\subsection*{A.5 Pop~III Star Formation Rate}

The Pop~III SFR becomes
\begin{equation}
    \dot{\rho}_{\star,\mathrm{III,RDCC}}
    = \dot{\rho}_{\star,\mathrm{III},\Lambda\mathrm{CDM}}
      \left[
        1 - \chi_{\mathrm{III}}
        \left(\frac{M}{M_{\mathrm{rel}}}\right)^{-\alpha}
      \right],
\end{equation}
with $\chi_{\mathrm{III}} \sim 0.2$.

This approximation matches the full RDCC Pop~III SFR to within $10\%$.

% ---------------------------------------------------------
\subsection*{A.6 Pop~II Star Formation Rate}

The Pop~II SFR becomes
\begin{equation}
    \dot{\rho}_{\star,\mathrm{II,RDCC}}
    = \dot{\rho}_{\star,\mathrm{II},\Lambda\mathrm{CDM}}
      \left[
        1 - \chi_{\mathrm{II}}
        \left(\frac{M}{M_{\mathrm{rel}}}\right)^{-\alpha}
      \right],
\end{equation}
with $\chi_{\mathrm{II}} \sim 0.1$.

This reproduces the RDCC Pop~II suppression to within $5$–$8\%$.

% ---------------------------------------------------------
\subsection*{A.7 Stellar Mass–Halo Mass Relation}

The stellar mass becomes
\begin{equation}
    M_{\star,\mathrm{RDCC}}
    = M_{\star}
      \left[
        1 - \chi_{\star}
        \left(\frac{M}{M_{\mathrm{rel}}}\right)^{-\alpha}
      \right],
\end{equation}
with $\chi_{\star} \sim 0.1$.

This matches the RDCC stellar mass suppression to within $10\%$.

% ---------------------------------------------------------
\subsection*{A.8 UV Luminosity Function}

The RDCC UV luminosity function is approximated by
\begin{equation}
    \phi_{\mathrm{RDCC}}(L)
    = \phi(L)
      \left[
        1 - \chi_{\mathrm{UV}}
        \left(\frac{L}{L_{\mathrm{rel}}}\right)^{-\alpha}
      \right],
\end{equation}
with $\chi_{\mathrm{UV}} \sim 0.1$.

This captures the RDCC suppression of the faint-end slope.

% ---------------------------------------------------------
\subsection*{A.9 Summary}

The analytical approximations derived in this appendix provide:

\begin{itemize}
    \item compact expressions for RDCC-modified early-galaxy processes,
    \item fast semi-analytic tools for parameter inference and MCMC scans,
    \item and physically transparent insight into the early-galaxy signatures of RDCC.
\end{itemize}

These results complement the numerical analysis of the main text and provide a practical 
framework for modeling early galaxy formation in RDCC.

% ---------------------------------------------------------
% APPENDIX B — IMPLEMENTATION OF RDCC–EARLY-GALAXY PHYSICS IN CLASS
% ---------------------------------------------------------
\section*{Appendix B: Implementation of RDCC–Early-Galaxy Physics in CLASS}
\addcontentsline{toc}{section}{Appendix B: Implementation of RDCC–Early-Galaxy Physics in CLASS}

This appendix describes the implementation of the RDCC early-galaxy formation model 
(RDCC–EGF) in the CLASS Boltzmann code. The modifications extend the RDCC linear and 
non-linear framework developed in Companion~XVI–XIX and incorporate the cooling, fragmentation, 
and star formation processes derived in Companion~XXII.

The goal is to compute the RDCC-modified halo abundances, cooling thresholds, star formation 
rates, stellar masses, and UV luminosities, and to make these quantities available to CLASS 
and external post-processing tools.

% ---------------------------------------------------------
\subsection*{B.1 Required CLASS Modules}

The RDCC–EGF implementation requires modifications in the following CLASS modules:

\begin{itemize}
    \item \texttt{background.c} — access to $k_{\mathrm{rel}}(a)$ and RDCC growth functions,
    \item \texttt{thermodynamics.c} — gas temperatures and cooling functions,
    \item \texttt{reionization.c} — consistency with RDCC ionization history,
    \item \texttt{common.h} — new RDCC–EGF data structures,
    \item \texttt{input.c} — new user parameters for early-galaxy physics,
    \item \texttt{output.c} — UV luminosity function and stellar mass outputs.
\end{itemize}

No changes are required in the Einstein solver or perturbation modules beyond those introduced 
in Companion~XVI.

% ---------------------------------------------------------
\subsection*{B.2 Data Structures}

CLASS must store the following RDCC–EGF quantities:

\begin{itemize}
    \item relaxation scale $k_{\mathrm{rel}}$,
    \item RDCC HMF suppression parameters $(\chi_{\mathrm{HMF}}, \alpha)$,
    \item cooling suppression parameters $(\chi_{\mathrm{cool}}, \chi_{\mathrm{H_2}})$,
    \item fragmentation and Jeans mass parameters $(\chi_{\mathrm{J}})$,
    \item Pop~III and Pop~II SFR suppression parameters $(\chi_{\mathrm{III}}, \chi_{\mathrm{II}})$,
    \item stellar mass suppression parameter $\chi_{\star}$,
    \item UV luminosity suppression parameter $\chi_{\mathrm{UV}}$.
\end{itemize}

These are stored in a dedicated \texttt{rdcc\_egf} structure in \texttt{common.h}.

% ---------------------------------------------------------
\subsection*{B.3 Input Parameters}

The following new input parameters are added to \texttt{input.c}:

\begin{itemize}
    \item \texttt{rdcc\_chi\_HMF},
    \item \texttt{rdcc\_chi\_cool},
    \item \texttt{rdcc\_chi\_H2},
    \item \texttt{rdcc\_chi\_J},
    \item \texttt{rdcc\_chi\_III},
    \item \texttt{rdcc\_chi\_II},
    \item \texttt{rdcc\_chi\_star},
    \item \texttt{rdcc\_chi\_UV},
    \item \texttt{k\_rel}.
\end{itemize}

Default values match the analytical approximations of Appendix~A.

% ---------------------------------------------------------
\subsection*{B.4 RDCC Halo Mass Function in CLASS}

The RDCC-modified halo mass function is computed as
\begin{equation}
    \frac{dn_{\mathrm{RDCC}}}{dM}
    = \frac{dn_{\Lambda\mathrm{CDM}}}{dM}
      \left[
        1 - \chi_{\mathrm{HMF}}
        \left(\frac{M}{M_{\mathrm{rel}}}\right)^{-\alpha}
      \right].
\end{equation}

CLASS must:

\begin{itemize}
    \item compute the standard Sheth–Tormen HMF,
    \item apply the RDCC suppression factor,
    \item tabulate $dn/dM$ on a logarithmic mass grid,
    \item store the result in \texttt{background->rdcc\_egf}.
\end{itemize}

% ---------------------------------------------------------
\subsection*{B.5 Cooling Thresholds and Gas Temperatures}

The molecular cooling threshold becomes
\begin{equation}
    M_{\mathrm{cool,RDCC}}
    = M_{\mathrm{cool}}
      \left[
        1 + \chi_{\mathrm{cool}}
        \left(\frac{M}{M_{\mathrm{rel}}}\right)^{-\alpha}
      \right].
\end{equation}

CLASS must:

\begin{itemize}
    \item compute $T_{\mathrm{vir}}(M,z)$,
    \item compute $\mathrm{H}_2$ cooling rates using RDCC-modified densities,
    \item evaluate the RDCC cooling threshold,
    \item store the result in \texttt{thermo->rdcc\_cool}.
\end{itemize}

% ---------------------------------------------------------
\subsection*{B.6 Pop~III and Pop~II Star Formation Rates}

The Pop~III SFR is
\begin{equation}
    \dot{\rho}_{\star,\mathrm{III,RDCC}}
    = \dot{\rho}_{\star,\mathrm{III}}
      \left[
        1 - \chi_{\mathrm{III}}
        \left(\frac{M}{M_{\mathrm{rel}}}\right)^{-\alpha}
      \right].
\end{equation}

The Pop~II SFR is
\begin{equation}
    \dot{\rho}_{\star,\mathrm{II,RDCC}}
    = \dot{\rho}_{\star,\mathrm{II}}
      \left[
        1 - \chi_{\mathrm{II}}
        \left(\frac{M}{M_{\mathrm{rel}}}\right)^{-\alpha}
      \right].
\end{equation}

CLASS must:

\begin{itemize}
    \item compute SFRs on the RDCC-modified HMF grid,
    \item integrate over halo mass,
    \item store the results in \texttt{background->rdcc\_sfr}.
\end{itemize}

% ---------------------------------------------------------
\subsection*{B.7 Stellar Mass and UV Luminosity}

The stellar mass becomes
\begin{equation}
    M_{\star,\mathrm{RDCC}}
    = M_{\star}
      \left[
        1 - \chi_{\star}
        \left(\frac{M}{M_{\mathrm{rel}}}\right)^{-\alpha}
      \right].
\end{equation}

The UV luminosity function becomes
\begin{equation}
    \phi_{\mathrm{RDCC}}(L)
    = \phi(L)
      \left[
        1 - \chi_{\mathrm{UV}}
        \left(\frac{L}{L_{\mathrm{rel}}}\right)^{-\alpha}
      \right].
\end{equation}

CLASS must:

\begin{itemize}
    \item compute $M_{\star}(M,z)$,
    \item convert SFR to UV luminosity,
    \item tabulate the UV-LF on a luminosity grid,
    \item output the UV-LF via \texttt{output.c}.
\end{itemize}

% ---------------------------------------------------------
\subsection*{B.8 Numerical Stability}

To ensure numerical stability:

\begin{itemize}
    \item use logarithmic mass grids for HMF and SFR,
    \item spline-interpolate all RDCC–EGF quantities,
    \item precompute scale-dependent factors involving $k_{\mathrm{rel}}$,
    \item avoid evaluating suppression factors at extremely small masses,
    \item ensure smooth transitions between Pop~III and Pop~II regimes.
\end{itemize}

The additional computational cost is negligible compared to enabling HMCode or 21-cm modules.

% ---------------------------------------------------------
\subsection*{B.9 Summary}

This appendix provides a complete implementation strategy for the RDCC early-galaxy formation 
model in CLASS. The modifications are minimal, the numerical cost is low, and the resulting 
stellar masses, star formation rates, and UV luminosity functions are fully consistent with 
the analytical and semi-analytic results derived in the main text.

% ---------------------------------------------------------
% APPENDIX C — COMPARISON WITH JWST AND HIGH-z GALAXY OBSERVATIONS
% ---------------------------------------------------------
\section*{Appendix C: Comparison with JWST and High-$z$ Galaxy Observations}
\addcontentsline{toc}{section}{Appendix C: Comparison with JWST and High-\texorpdfstring{$z$}{z} Galaxy Observations}

This appendix compares the predictions of the Relaxation-Driven Cyclic Cosmology (RDCC) for 
early galaxy formation with current observational constraints from JWST, HST legacy fields, 
quasar damping wings, and indirect 21-cm limits. These observations probe the stellar mass 
buildup, UV luminosity function (UV-LF), star formation rates (SFRs), clustering, and 
metallicity of galaxies at $z > 6$.

% ---------------------------------------------------------
\subsection*{C.1 UV Luminosity Function}

JWST observations show:

\begin{itemize}
    \item a suppressed faint-end slope at $z > 10$,
    \item a slower evolution of the UV-LF than predicted by $\Lambda$CDM,
    \item fewer low-luminosity galaxies than expected.
\end{itemize}

RDCC predicts:

\begin{itemize}
    \item reduced star formation efficiency in low-mass halos,
    \item delayed cooling and gas accretion,
    \item suppressed faint-end UV-LF due to the RDCC HMF suppression.
\end{itemize}

\paragraph{Key agreement.}
RDCC naturally reproduces the shallow faint-end slope and slow evolution of the UV-LF observed 
by JWST without requiring additional feedback or tuning.

% ---------------------------------------------------------
\subsection*{C.2 Stellar Masses and Star Formation Rates}

JWST reveals:

\begin{itemize}
    \item lower stellar masses at $z > 10$ than predicted by $\Lambda$CDM,
    \item reduced specific star formation rates (sSFRs),
    \item smoother stellar mass buildup.
\end{itemize}

RDCC predicts:

\begin{itemize}
    \item reduced baryon fractions in low-mass halos (Companion~XX),
    \item delayed Pop~III and Pop~II star formation (Section~4),
    \item slower stellar mass growth (Section~5).
\end{itemize}

\paragraph{Key agreement.}
RDCC matches the suppressed stellar masses and reduced sSFRs observed by JWST.

% ---------------------------------------------------------
\subsection*{C.3 Galaxy Clustering}

JWST and HST clustering analyses show:

\begin{itemize}
    \item enhanced clustering of galaxies at $z > 8$,
    \item bias values higher than expected for their stellar masses,
    \item strong large-scale correlations.
\end{itemize}

RDCC predicts:

\begin{itemize}
    \item enhanced galaxy bias due to suppressed $\sigma(M)$ (Section~5),
    \item stronger clustering for halos near the relaxation scale $k_{\mathrm{rel}}$,
    \item reduced abundance of low-mass halos, increasing the relative weight of massive halos.
\end{itemize}

\paragraph{Key agreement.}
RDCC naturally explains the enhanced clustering of early galaxies without invoking exotic 
feedback or modified initial conditions.

% ---------------------------------------------------------
\subsection*{C.4 Quasar Damping Wings and Metallicity}

High-redshift quasars provide constraints on:

\begin{itemize}
    \item the neutral fraction of the IGM,
    \item the metallicity of early galaxies,
    \item the distribution of ionized bubbles.
\end{itemize}

Observations show:

\begin{itemize}
    \item moderate metallicities at $z \sim 7$,
    \item large ionized regions around quasars,
    \item significant sightline variation.
\end{itemize}

RDCC predicts:

\begin{itemize}
    \item delayed but extended metal enrichment (Section~4),
    \item larger ionized regions due to enhanced escape fractions (Companion~XXI),
    \item smoother damping wings due to reduced small-scale structure.
\end{itemize}

\paragraph{Key agreement.}
RDCC matches the observed metallicities and ionized bubble sizes without requiring extreme 
Pop~III star formation.

% ---------------------------------------------------------
\subsection*{C.5 Indirect 21-cm Constraints}

Although no detection exists yet, 21-cm upper limits imply:

\begin{itemize}
    \item no strong early heating,
    \item no deep absorption trough at $z \sim 15$--$20$,
    \item delayed coupling of the spin temperature.
\end{itemize}

RDCC predicts:

\begin{itemize}
    \item reduced early star formation (Section~4),
    \item delayed Pop~III heating,
    \item weaker global 21-cm absorption (Companion~XXI).
\end{itemize}

\paragraph{Key agreement.}
RDCC is consistent with all current 21-cm upper limits.

% ---------------------------------------------------------
\subsection*{C.6 Summary Table}

\begin{table}[H]
\centering
\begin{tabular}{lccc}
\hline
\textbf{Observable} & \textbf{Data} & \textbf{RDCC Prediction} & \textbf{Agreement} \\
\hline
UV-LF faint-end slope & shallow & suppressed faint-end & strong \\
Stellar masses & low at $z>10$ & reduced buildup & strong \\
sSFR & low & reduced SFE & strong \\
Clustering & enhanced & higher bias & strong \\
Metallicity & moderate & delayed enrichment & strong \\
21-cm limits & weak heating & delayed Pop~III & strong \\
\hline
\end{tabular}
\caption{Comparison of RDCC predictions with JWST and high-$z$ galaxy observations.}
\label{tab:comparisonCXXII}
\end{table}

% ---------------------------------------------------------
\subsection*{C.7 Concluding Remarks}

RDCC differs fundamentally from $\Lambda$CDM in that:

\begin{itemize}
    \item early galaxies form later and more smoothly,
    \item stellar masses and SFRs are reduced at $z > 10$,
    \item the UV-LF faint-end slope is suppressed,
    \item early galaxies are more strongly clustered,
    \item metal enrichment is delayed,
    \item and the relaxation scale $k_{\mathrm{rel}}$ imprints a characteristic signature.
\end{itemize}

These differences provide clear observational signatures for current and upcoming surveys, 
especially JWST, Roman, and 21-cm experiments.

\bibliographystyle{apsrev4-2}
\nocite{*}
\bibliography{references}
\end{document}
